%%%%% Définition paramètres documents

\documentclass[a4paper, 12pt]{report} % Initilisation document

% Paramètres texte

\usepackage[utf8]{inputenc} % Français avec caractères spéciaux
\usepackage[french]{babel} % Français pour la biographie
\usepackage[T1]{fontenc} % Polices d'écriture

% Mise en page

\usepackage{bookmark}
\usepackage[left=2cm, right=2cm, top=2cm, bottom=2cm]{geometry} % Dimensions marges
\usepackage{indentfirst} 
\usepackage{underscore}
\usepackage{nopageno} % Supprimer la pagination d'une page
\usepackage{float, caption}	% Positionnement et légende des images
\usepackage{fancyhdr} % Marges
\usepackage{calc, ifthen, xspace} % Redéfinition des espaces et distances

% Paramètres mathématiques

\usepackage{amsmath, amsthm, amssymb, amsfonts} % Polices et symboles mathématiques sophistiqués

% Paramètres graphiques

\usepackage{graphicx} % Inclusion images
\usepackage{epic,eepic}	% Positionnement image de garde
\usepackage{wrapfig} % Images sur le coté dans le texte

% Autres

\usepackage{hyperref} % Liens croisés fichier PDF
\usepackage{color} % Couleur codes informatiques
\usepackage{lastpage} % Références aux pages
\usepackage{endnotes} % Notes de fin

% Identation automatique

\setlength{\parindent}{0pt} % Supression des alinéas automatiques par défaut, plus manipulable

% Paramètres hyperliens

\hypersetup{
% dvips,
% backref=true,
% pagebackref=true, % Bibliographies
% hyperindex=true, % Index.
colorlinks = true, % Couleur hyperliens
breaklinks = true, % Retour à la ligne dans les hyperliens trop longs
urlcolor = blue, % Couleur des hyperliens
linkcolor = black, % Couleur des liens internes
% bookmarks=blue, % Signets pour Acrobat
bookmarksopen = true}

\makeatletter
\g@addto@macro{\UrlBreaks}{\do\-\do\.}
\makeatother

% Paramètres de différentes couleurs avec {nom_couleur}{type_encodage}{proportions}

\definecolor{hellgelb}{rgb}{1,1,0.8}
\definecolor{colKeys}{rgb}{0,0,1}
\definecolor{colIdentifier}{rgb}{0,0,0}
\definecolor{colComments}{rgb}{0,0.5,0}
\definecolor{colString}{rgb}{0.62,0.12,0.94}
\definecolor{INSA_GM}{cmyk}{0.6,0,0,0}
\definecolor{INSA_GRIS}{cmyk}{0.7,0.6,0.5,0.3}
\definecolor{INSA_BLEU}{cmyk}{1,0.9,0.1,0}

% Définition de commandes compliquées pour déclarer un nom de projet

\newcommand{\insertrefproj}[1]{}
\newcommand{\refproj}[1]{\renewcommand{\insertrefproj}{\textbf{\color{INSA_GRIS}#1}}}

% Style de page

\pagestyle{fancy}

\fancypagestyle{courant}{
	\fancyhf{}
    \setlength{\headheight}{27pt}
    \fancyhead[L]{\raisebox{-2mm}{\includegraphics[width=30mm]{Images/Logo INSA.png}}}
\fancyhead[C]{}
\fancyhead[R]{\color{INSA_GRIS}\thepage}%\scshape\leftmark}
\fancyfoot[L]{\insertrefproj}
\fancyfoot[R]{}
%\fancyfoot[CE,CO]{\color{INSA_GRIS}\thepage}
\renewcommand{\headrulewidth}{0pt}
\renewcommand{\footrulewidth}{0.2pt}
}

\fancypagestyle{special}{%
  \pagestyle{courant}
  \fancyfoot{}
  %\addtolength{\footskip}{-20pt}
%  \setlength{\footskip}{0pt}
%  \fancyfoot[C]{}
\renewcommand{\footrulewidth}{0pt}
}

\fancypagestyle{plain}{%
  \fancyhf{}%
  \pagestyle{courant}
}

% 

\newcommand{\Scilab}{
	\lstset{
	language=Scilab,
	float=hbp,
	basicstyle=\ttfamily\small,
	identifierstyle=\color{colIdentifier},
	keywordstyle=\bf \color{colKeys},
	stringstyle=\color{colString},
	commentstyle=\color{colComments},
	columns=flexible,
	tabsize=5,
	frame=single,
	%frame=shadowbox,
	rulesepcolor=\color[gray]{0.5},
	extendedchars=true,
	showspaces=false,
	showstringspaces=false,
	numbers=left,
	stepnumber=5,
	firstnumber=1,
	numberstyle=\tiny,
	breaklines=true,
	%backgroundcolor=\color{hellgelb},
	captionpos=b,
	}
}

% Paramètres projet

\title{Projet de P6}
\refproj{STPI/P6/2024 - 39}
\author{}
\date{}

%%%%% Début d'édition du document

\begin{document}

%%% Page de présentation

\pagenumbering{Alph} % Numérotation avec lettre majuscule
\thispagestyle{empty} % Type page vide intiallement

\vspace{4cm}

\begin{picture}(0,0) % Hauteur, Largeur
	
\put(80,-465){\includegraphics[scale=1]{Images/Image Garde.png}} % Image de garde centrée horizontalement
\put(0,-20){\includegraphics[width=0.4\textwidth]{Images/Logo INSA.png}} % Logo INSA haut-gauche
\put(240,0){{\begin{minipage}{12cm}\centering \Large % Caractéristiques projet haut droit
	\textbf{Projet de Physique P6} \\ 
	\textbf{STPI/P6/2024 - 39}\end{minipage}}}
\put(-20,-150){{\begin{minipage}{\textwidth}\centering \Huge % Titre princpal rapport
	\textbf{Modélisation de l'impact des gaz à effet de serre}\end{minipage}}}
% Quelques changements pour illustrer
\newsavebox{\noms}
\savebox{\noms}(300, 300)[tl]{
\put(0,98){\color{INSA_GRIS}{\begin{minipage}{6cm} 
	\textbf{Étudiants :} \end{minipage}}} % Étudiants par ordre alphabétique de nom
\put(0,60){\color{INSA_GRIS}{\begin{minipage}{6cm}
	Romane LANÈRES \\ Anouk PETITGAS \\ Arthur SARRAU \end{minipage}}}
\put(185,60){\color{INSA_GRIS}{\begin{minipage}{6cm}
	Clara MÉLINE \\ Tom PHILIPPE \\ Nina ZEDDOUN \end{minipage}}}

\put(0,){\color{INSA_GRIS}\begin{minipage}{9cm}   
	\textbf{Enseignant-responsable du projet :} \\ Samuel PAILLAT \end{minipage}}}

\put(100,-850){\usebox{\noms}}

\end{picture}

%%% Page complètement vierge

\newpage
\pagenumbering{arabic}
\setcounter{page}{1}
\thispagestyle{empty}
\null

%%% Page descriptions cractéristiques principales projet

\newpage
\pagestyle{special}

\chapter*{Fiche projet} % * Garde la même mise en page qu'un chapitre mais ne le numérote pas
\addcontentsline{toc}{chapter}{Fiche projet} % Ajout dans le sommaire, possiblement renommé comme souhaité

\textbf{Date de remise du rapport :} 13/06/2024 \vspace{\baselineskip}
% \vspace{\baselineskip} permet de sauter des lignes dans le code (pour la lisibilité)
% en évitant tout message d'erreur de la forme "Underfull \hbox (badness 10000)"

\textbf{Référence du projet:} STPI/P6/2024 - 39 \vspace{\baselineskip}

\textbf{Intitulé du projet :} Modélisation de l'impact des gaz à effet de serre \vspace{\baselineskip}

\textbf{Type de projet :} Modélisation numérique \vspace{\baselineskip}

\textbf{Objectifs du projet :} \vspace{\baselineskip} 

Ce projet a pour vocation de développer un modèle numérique
décrivant les mécanismes de régulation thermique de 
l'atmosphère terrestre. Premièrement, nous allons réaliser
une étude à travers quelques calculs en ordre de grandeurs afin
d'identifier les différents paramètres d'influence de notre modèle.
L'objectif est de mettre au point un modèle qui soit modulable
en termes de paramètres et de complexité. Cette dernière est 
ajustée par le traitement de différents gaz à effet de serre 
et de divers modèles de température et de pression. 
L'ensemble des modélisations numériques aspire à proposer 
une version améliorée du travail réalisé par le 
vulgarisateur scientifique David Louapre. À travers ce projet, 
nous cherchons à reproduire au mieux le 
phénomène d'équilibre thermique de la Terre en considérant 
les flux principaux et en quantifiant précisément l'impact
des gaz à effet de serre dans le système atmosphérique. \vspace{\baselineskip}

\textbf{Mots-clefs du projet :} Transfert thermique, Corps noir, Effet de serre, Bilan radiatif 
\vspace{\baselineskip} 

\vfill

% Coordonnées INSA Rouen par défaut

\begin{center}
	\scshape Institut National des Sciences Appliquées de Rouen \\
	Département Sciences et Techniques Pour l'Ingénieur \\
	685 Avenue de l'Université BP 08- 76801 Saint-Etienne-du-Rouvray \\ Tél : 33 2 32 95 66 21 - Fax : 33 2 32 95 66 31
\end{center}

%%% Remerciements

\newpage
\chapter*{Remerciements} 
\addcontentsline{toc}{chapter}{Remerciements}

\setlength{\parindent}{30pt}

\indent Nous souhaiterions avant tout remercier Samuel PAILLAT, notre
coordinateur et enseignant encadrant, qui a su nous aiguiller
et nous apporter son aide et ses connaissances tout au longs
de ce projet. \vspace{\baselineskip}

	Aussi nous aimerions nous porter reconnaissants 
envers le physicien médaillé de la médiation scientifique
du CNRS: David Louapre. Ces travaux sur l'effet de serre 
ont été une grande source d'inspiration en tant que point
de départ pour notre recherche \endnote{DAVID LOUAPRE, Article et vidéo,
\textit{"La saturation de l’effet de serre (ou pas)"}, 06/10/2023:
\\ \url{https://scienceetonnante.com/2023/10/06/la-saturation-de-leffet-de-serre/}}. 
\vspace{\baselineskip}

\indent Nous saluons également les divers organismes ayant 
mis à disposition des bases de données en open source. Nous
en avons extrait des paramètres physiques sans lesqueles 
nos modèles auraient manqué de fiabilité. \vspace{\baselineskip}

\indent Enfin, nous aimerions remercier l'organisation de 
l'INSA plus généralement pour nous avoir implémenté l'EC de P6
qui nous a permis de réaliser notre premier projet
scientifique appliqué dans des conditions de collaboration
collective proches de celle milieu de l'ingénierie. 
En effet depuis le début de ce semestre nous avons tous commencé à 
nous pré-spécialiser dans un domaine de l'ingénierie. À travers 
ce projet chacun est amené à mettre en commun ses savoirs et 
compétences acquis pendant ces derniers mois, ce qui s'avère est 
enrichissant pour tout le groupe.

\vfill

%%% Table des matières 

\newpage
\pagestyle{courant} 
	\setcounter{tocdepth}{2} % Règle la profondeur maximale de la table des matières
	\tableofcontents % Affichage table des matières

%%% Notations et Acronymes

\newpage
\chapter*{Notations} 
\addcontentsline{toc}{chapter}{Notations}

\begin{description}
	\subsubsection*{Acronymes}
	\item[GES:] Gaz à Effet de Serre
    \item[CN:] Corps Noir
    
	\subsubsection*{Grandeurs \footnote{En unités usuelles}}
    \item[\boldmath{$\Phi$}:] Flux ou débit \footnote{Dans ce rapport 
	il décrira exclusivement la puissance rayonnée par unité de surface} ($W.m^{-2}$)
    \item[\boldmath{$M^{0}$}:] Émittance totale d'un corps noir, ($W.m^{-2}$)
    \item[\boldmath{$M^{0}_{\lambda,T}$}:] Émittance d'un corps noir ($W.m^{-3}$)
    \item[\boldmath{$\epsilon$}:] Émissivité (sans unité)
    \item[\boldmath{$\tau_{CO_2}$}:] Coefficient de transmission du $CO_2$ (sans unité)
    \item[\boldmath{$\lambda$}:] longueur d'onde ($m$)
    \item[\boldmath{$P$}:] Pression ($Pa$)
    \item[\boldmath{$V$}:] Volume ($m^3$)
	\item[\boldmath{$n$}:] Quantité de matière ($mol$)
	\item[\boldmath{$T$}:] Température ($K$)
 	\item[\boldmath{$M$}:] Masse molaire ($kg.mol^{-1}$)
	
	\subsubsection*{Constantes \footnote{En Unités du Système International (USI) à 3 CS près}}
	\item[\boldmath{$R$}:] Constante universelle des gaz parfaits = $8,314$ $kg.m^2.mol^{-1}.K^{-1}.s^{-2}$
	\item[\boldmath{$\sigma$}:] Constante de Stefan-Boltzmann = $5,67.10^{-8}$ $W.m^{-2}.K^{4}$
	\item[\boldmath{$h$}:] Constante de Plank = $6,63.10^{-34}$ $kg.m^2.s^{-1}$
 	\item[\boldmath{$k_B$}:] Constante de Boltzmann = $1,38^{-23}$ $kg.m^2.s^{-2}.K^{-1}$
  	\item[\boldmath{$g$}:] Constante gravitationelle terrestre = $9,81$ $m.s^{-2}$
   	\item[\boldmath{$c_0$}:] Célérité de la lumière = $2,99^8$ $m.s^{-1}$
    \item[\boldmath{$C_1$}:] Constante de Planck 1 = $1,19.10^{-16}$ $kg.s^{-3}.m^{-4}$
    \item[\boldmath{$C_2$}:] Constante de Planck 2 = $1,44.10^{-2}$ $m.K$
    \item[\boldmath{$T_s$}:] Température moyenne du Soleil = $5,772.10^3$ $K$
    \item[\boldmath{$T_s$}:] Température moyenne de la Terre = $288,15$ $K$
	\item[\boldmath{$\bar{a_T}$}:] Albédo terrestre moyen = $0,31$
\end{description}

%%% Introduction

\newpage
\chapter*{Introduction}				
\addcontentsline{toc}{chapter}{Introduction}

\indent Le réchauffement climatique est sans aucun doute la plus grande problématique 
scientifique du siècle actuel. Elle regroupe une multitude d'enjeux mêlant environnement,
climat, industries et sociétés humaines. Cependant s'il y a un point en particulier qui 
fait parler de lui, il s'agit bien de l'effet de serre. \vspace{\baselineskip}

\indent Cependant, il semblerait que ce phénomène physique, vraisemblablement assez simple 
en apparence, s'avère bien plus complexe en réalité. Si l'on se penche dans un premier temps sur 
l'éthymologie de ce terme, inventé par le scientifique français Joseph Fourier
\endnote{JOSEPH FOURIER, Article, 
\textit{"Mémoire sur les températures du globe terrestre et des espaces planétaires}, 
Mémoires de l'Académie royale des sciences de l'Institut de France, vol.7,  p.569-604, 1827:
\\ \url{https://gallica.bnf.fr/ark:/12148/bpt6k32227/f814.item}},
on recontre déjà une première confusion notable. En effet l'effet de serre 
tel qu'il opère dans les couches atmosphériques est majoritairement lié au jeu d'absorption, d'émission 
et de transmission de rayonnements. Tandis que les serres misent plutôt sur l'utilisation d'une enceinte 
fermée hermétique qui supprime l'effet de convection de l'air, ce qui a pour conséquence une élévation 
de température. En bref, il faut garder en tête que cette analogie comporte des limites. \vspace{\baselineskip}

\indent De ce fait, c'est dans ce cadre que s'inscrit notre projet, qui a pour but de quantifier et d'étudier
l'impact des gaz à effet de serre avec des outils de deuxième année de cursus ingénieur. L'objectif est 
de construire petit à petit un modèle, avec de plus en plus de paramètres, et qui tende à se rapprocher
d'une descritpion fidèle de l'impact des gaz à effet de serre. \vspace{\baselineskip}

\indent Tout d'abord nous réaliserons quelques calculs intiaux de flux sans aucun GES.
Ensuite nous introduirons le dioxyde de carbone $CO_2$ et étudierons son impact avec une température de
système uniforme. Puis nous tenterons de considérer d'autres GES et de faire évoluer leur caractéristiques
en fonction de modèles atmosphériques de température, altitude et pression. \vspace{\baselineskip}

\indent Dans l'ensemble, ce projet se base sur la physique des transferts thermiques, 
mais fait également appel à des compétences en informatique et mathématiques.
En effet notre modèle se fonde sur plusieurs paramètres qu'il sera nécessaire
d'ajuster selon nos besoins, ce qui requiert d'implémanter tous ces calculs
numériquement. \vspace{\baselineskip}

\indent Pour ce qui est des outils, nous avons codé l'ensemble de nos algorithmes en Python 3.12 et avons misé
la popularité de ce language pour utiliser des bibliothèques de physique et de mathématiques libres d'accès
permettant d'avoir accès à des outils absolument fondamentaux: calcul intégral numérique, 
base de donnée de taux d'absorptions de gaz à effet de serre, graphiques... Sinon, nous avons utilisé du \LaTeX
pour la rédaction de ce rapport, qui est un logiciel parfaitement adapté aux études dans le milieu scientifique.
La prise en main s'est faite progressivement pour chaque membre du groupe, mais au final, nous sommes tous parvenus
à maîtriser les bases fondamentales et rudiments de cet outil de traitement de texte à première vue bien complexe.

%%% Méthodologie, organisation du travail

\chapter{Méthodologie, organisation du travail}

\begin{itemize}
\item Description de l'organisation adoptée pour le déroulement du travail
\item Organigramme des tâches réalisées et des étudiants concernés
\end{itemize}

%%% Travail réalisé et résultats

\chapter{Travail réalisé et résultats}

\section{L'effet de serre}

\subsection{Explication du phénomène de l'effet de serre}

	L'effet de serre est un processus naturel qui se produit 
avec le rayonnement entre la Terre et le Soleil. Tout d'abord, 
le Soleil émet un rayonnement vers la Terre, l'atmosphère 
laisse passer une partie de ce rayonnement solaire. Réchauffé, 
le corps terrestre, étant un corps noir, émet un flux rayonné 
identique à celui reçu par soleil. Ce flux est alors renvoyé 
vers l'atmosphère et une partie est absorbée par les gaz 
présents, appelés gaz à effet de serre. Le reste du flux 
est soit envoyé dans l'espace, soit renvoyé vers la Terre. 
De plus, nous avons pu constater que ce phénomène n'existe 
que lorsque que la température varie dans l'atmosphère en 
fonction de l'altitude. Ainsi si la température était uniforme 
dans l'atmosphère ce phénomène n'existerait pas. \vspace{\baselineskip}

	C'est l'absorbtion du rayonnement par les GES qui va 
provoqué un rechauffement terrestre. Ainsi, ce phénomène 
créé un réchauffement global de la Terre permettant à la 
température moyenne de s'éléver à 15°C au lieu de -18°C 
(sans GES). 

\subsection{Définition des GES}
Gaz à effet de serre est gaz présent dans l'atmosphère qui 
absorbe une partie des rayonnements (rayons infrarouges) 
reçue par le soleil. Ces gaz sont d'origine naturelle 
(vapeur d'eau) et/ou anthropique 
(issues des activités humaines) comme par exemple le dioxyde 
de carbonne ($CO_2$), le méthane ($CH_4$), le protoxyde d'azote
($N_2O$) et les gaz fluorés (HFC) \endnote{JEAN-MARC JANCOVICI, Article, 
\textit{"Quels sont les gaz à effet de serre ?"}, 01/08/2007:
\\ \url{https://jancovici.com/changement-climatique/gaz-a-effet-de-serre-et-cycle-du-carbone/quels
-sont-les-gaz-a-effet-de-serre-quels-sont-leurs-contribution-a-leffet-de-serre/}}. \vspace{\baselineskip}

% mettre schema flux radiatif 

\section{Bilan radiatif terrestre}
\subsection{Définitions du corps noir et de l'émittance}

\subsubsection{Le corps noir}
Un corps noir est un objet idéal qui permet d'évaluer le flux thermique maximum
que peut rayonner un corps, en fonction de sa température. Ainsi un corps noir 
va émettre autant de rayonnement qu'il en absorbe. On peut donc dire que son 
emmitance $ \epsilon$  est égale à 1. 

\subsubsection{L'émittance} 

La luminance d'un corps noir est évalué grace à la loi de Planck 
(en $W.m^{-3}$): \vspace{\baselineskip}

\begin{center}
$M^{0}_\lambda,_T= \frac{\pi C_1}{\lambda^5} \biggl(\exp(\frac{C_2}{\lambda T}) - 1\biggr)^{-1}$ 

$C_1$ = 1,19 $\cdot 10^{8}$ $W.m^{-2}.\mu.m^{-1}$\
$C_2$ = 14400 $\mu.m.K$ \\
\end{center}
Et ainsi l'exitance totale d'un coprs noir (en $W.m^{-2}$) est obtenue avec la formule suivante:

\hfil $M^{0}$ = $\int_{0}^{\infty} M^{0}_\lambda d\lambda$ = $\sigma T^{4}$ 

\par

\subsection{Bilan radiatif sans effet de serre} 

\subsubsection{A température constante}

\subsubsection{A température variable en fonction de l'altitude}

Le Soleil émet des rayonnements dans le domaine visible et 
dans le domaine infrarouge. C'est un corps noir, c'est à dire
qu'il absorbe tout le rayonnement qu'il reçoit et le réémet 
parfaitement. De plus , on peut dire que son émissivité est 
égale à 1 d'après les hypothèses du corps noir. Il suit donc
la loi de Planck et son spectre d'émission en fonction de 
la longueur d'onde est le suivant :

% mettre image du spectre du soleil

On peut alors calculer ce qu'il émet en rayonnant de 2 manières: \vspace{\baselineskip}

1. Avec le flux radiatif
\begin{center}
$\Phi_{ray}$ = $\int \epsilon \sigma T^{4} dS$    
\end{center}

2. Numériquement
	En effet évaluer le flux émit revient à calculer l'aire sous la courbe 
de son spectre d'émission. (voir Figure  "mettre nombre figure" et inserer 
le code pour faire l'intégrale sous la courbe) \vspace{\baselineskip}

	Ainsi avec ces deux méthodes, on trouve que le rayonnement émis par le Soleil 
est de 174 PW  soit $1,74.10^{17}$ W

\chapter*{Conclusion et perspectives}
\addcontentsline{toc}{chapter}{Conclusion et perspectives}
\begin{itemize}
\item Conclusions sur le travail réalisé
\item Conclusions sur l'apport personnel de cet E.C. projet
\item Perspectives pour la poursuite de ce projet
\end{itemize}

%%% Bibliographie

\newpage

\renewcommand{\notesname}{} % Laisse vide le titre de la section
\chapter*{Bibliographie}
\addcontentsline{toc}{chapter}{Bibliographie}
% \renewcommand{\theendnote}{[\arabic{endnote}]} % Redéfinition de l'affichage des numéros de notes de fin
% \renewcommand{\enoteformat}{[\arabic{endnote}]
% 							\hspace{1em} % Laisse un petit espace entre l'indice et la note
% 							\newline} % Sauter une ligne entre chaque note 
\makeatletter
\renewcommand{\enoteheading}{\par\vspace{1 em}}
\renewcommand{\theenmark}{\makebox[0.5 em][r]{\@theenmark}}
\renewcommand{\enoteformat}{\parindent = 2 em
  							\leftskip = 0.5 em
  							[\theenmark]\enspace\ignorespaces}							
\makeatother
\theendnotes

%%% Annexes

\newpage
\chapter*{Annexes}
\addcontentsline{toc}{chapter}{Annexes}
\renewcommand{\thesection}{\Alph{section}}  % Renomme les futures sections avec lettres majuscules
\setcounter{section}{0} % Réinitialisation numéro (ici lettres) des sections à 0 (respectivement A)

\section{Documentation technique}
\newpage

\section{Listings des programmes réalisés}
\newpage

\section{Schémas de montages, plans de conception...}
\newpage

\section{Propositions de sujets de projets (en lien ou pas avec le projet réalisé)}
\newpage

\section{Mettre du code en annexe}

\end{document}
